% P. Areias simple C# -- Numerical & Scientific Computing Edition -- traditional Main, no ref/in/out teaching examples
% pdflatex csharp_perfect.tex (run twice for stable layout)
\documentclass[9pt,a4paper]{extarticle}
\usepackage[T1]{fontenc}
\usepackage[utf8]{inputenc}
\usepackage{helvet}
\usepackage{courier}
\renewcommand\familydefault{\sfdefault}
\usepackage[a4paper,top=7mm,bottom=7mm,left=6mm,right=6mm]{geometry}
\setlength{\parindent}{0pt}\setlength{\parskip}{0pt}
\pagestyle{empty}

\usepackage[dvipsnames,svgnames,x11names]{xcolor}
\definecolor{clKw}  {HTML}{7a0000}
\definecolor{clTy}  {HTML}{0a3d6b}
\definecolor{clFn}  {HTML}{4a1a6b}
\definecolor{clSt}  {HTML}{1e4d2b}
\definecolor{clCm}  {HTML}{5a4e3a}
\definecolor{clNu}  {HTML}{7a3800}
\definecolor{clOp}  {HTML}{1a0d00}
\definecolor{acBlue}  {HTML}{0a3d6b}
\definecolor{acTeal}  {HTML}{0a4e45}
\definecolor{acGreen} {HTML}{1e4d2b}
\definecolor{acPurple}{HTML}{4a1a6b}
\definecolor{acRed}   {HTML}{7a0000}
\definecolor{acAmber} {HTML}{7a3800}
\definecolor{acNavy}  {HTML}{0e1f4a}
\definecolor{acOrange}{HTML}{8a3200}
\definecolor{bgBlue}  {HTML}{f4f9fd}
\definecolor{bgTeal}  {HTML}{f0faf8}
\definecolor{bgGreen} {HTML}{f1f8f3}
\definecolor{bgPurple}{HTML}{f8f4fc}
\definecolor{bgRed}   {HTML}{fdf4f4}
\definecolor{bgAmber} {HTML}{fdf8ee}
\definecolor{bgNavy}  {HTML}{f3f5fb}
\definecolor{bgOrange}{HTML}{fdf4ee}
\definecolor{bgCode}  {HTML}{fefcfa}
\definecolor{bgHdr}   {HTML}{f7f3ee}
\definecolor{bdCard}  {HTML}{9a9080}
\definecolor{secBg}   {HTML}{1a0d00}
\definecolor{newTag}  {HTML}{4a7c59}
\definecolor{trap}    {HTML}{a02020}

\usepackage{listings}
\lstdefinelanguage{CSharp}{
  sensitive=true,
  morekeywords=[1]{abstract,as,base,bool,break,byte,case,catch,char,
    checked,class,const,continue,decimal,default,delegate,do,double,
    else,enum,event,explicit,extension,extern,false,field,file,
    finally,fixed,float,for,foreach,goto,if,implicit,in,init,int,
    interface,internal,is,lock,long,namespace,new,not,null,object,
    operator,override,params,partial,private,protected,public,
    readonly,record,ref,return,sbyte,sealed,short,sizeof,stackalloc,
    static,string,struct,switch,this,throw,true,try,typeof,uint,ulong,
    unchecked,unsafe,ushort,using,virtual,void,volatile,when,where,
    while,with,nint,nuint,global,async,await,var,and,or,scoped,
    required,allows,yield,get,set,value},
  morekeywords=[2]{Action,Activity,Array,ArrayPool,BinaryPrimitives,
    BitOperations,BinaryReader,BinaryWriter,Buffer,CancellationToken,
    CancellationTokenSource,Channel,CollectionsMarshal,Complex,
    ConcurrentBag,ConcurrentDictionary,ConcurrentQueue,Console,
    CultureInfo,DateOnly,DateTime,DateTimeOffset,Dictionary,
    Environment,Enumerable,EventCounter,File,Func,GC,GeneratedRegex,
    Guid,HashCode,HashSet,IAsyncDisposable,IAsyncEnumerable,
    IComparable,IComparer,IDisposable,IEnumerable,IEquatable,
    IEqualityComparer,IList,INumber,INumberBase,IFloatingPoint,
    IFloatingPointIeee754,IReadOnlyCollection,IReadOnlyDictionary,
    IReadOnlyList,ISet,ISpanFormattable,ISpanParsable,Interlocked,
    JsonSerializer,JsonSerializerOptions,JsonSourceGenerationOptions,
    JsonSerializerContext,JsonNamingPolicy,JsonNumberHandling,List,
    Lock,Math,MathF,Matrix4x4,Memory,MemoryMarshal,MemoryPool,
    NativeMemory,NumberStyles,ParallelOptions,Path,Predicate,
    PriorityQueue,Queue,Quaternion,Random,ReadOnlyMemory,ReadOnlySpan,
    Regex,RegexOptions,SearchValues,Span,Stack,Stopwatch,StreamReader,
    StreamWriter,StringBuilder,StringComparer,StringSyntaxAttribute,
    Task,TaskCompletionSource,Tensor,TensorPrimitives,TensorSpan,
    ThreadLocal,TimeOnly,TimeProvider,TimeSpan,Type,Unsafe,Uri,
    ValueTask,Vector,Vector2,Vector3,Vector4,Vector128,Vector256,
    Vector512,CallingConvention,IntPtr,Directory,
    ArgumentNullException,ArgumentOutOfRangeException,
    InvalidOperationException,KeyNotFoundException,
    OperationCanceledException,Enum,SortedDictionary,SortedSet,
    ArraySegment,MemoryExtensions,HttpClient,EventHandler,
    StringSplitOptions,FileNotFoundException,IOException,
    ArithmeticException,ObjectDisposedException},
  morekeywords=[3]{Abs,Acos,Add,AddRange,AddScaled,
    AllocateUninitializedArray,AllocateArray,All,Any,Asin,AsBytes,
    AsMemory,AsPointer,AsRef,AsSpan,Atan,Atan2,Average,Cancel,
    CancelAfter,Cast,Cbrt,Chunk,Clamp,Clear,Clone,CompareExchange,
    CompareTo,Concat,Contains,ContainsKey,ConvertChecked,CopyTo,Cos,
    Cosh,CosineSimilarity,Count,Create,CreateChecked,
    CreateSaturating,CreateScope,CreateSpan,CreateTruncating,
    CreateBounded,CreateUnbounded,Deserialize,Dequeue,DequeueMin,
    DistinctBy,Distance,Divide,Dot,EnsureCapacity,Enqueue,EnterScope,
    Equals,Exists,Exp,ExceptWith,Exchange,Fill,Find,FindIndex,First,
    FirstOrDefault,Floor,ForEach,FromResult,FromSeconds,FusedMultiplyAdd,
    GetCustomAttribute,GetHashCode,GetMethod,GetProperties,GetReference,
    GetValueRefOrAddDefault,GroupBy,Hypot,Increment,IndexOf,IndexOfAny,
    IndexOfAnyExcept,Insert,IntersectWith,Invoke,IsEvenInteger,IsFinite,
    IsHardwareAccelerated,IsInfinity,IsNaN,IsNegative,IsNormal,
    IsPositive,IsSubnormal,IsSubsetOf,IsSupported,IsZero,Length,
    LoadUnsafe,Log,Log10,Log2,LogP1,Match,Matches,Max,Min,Multiply,
    MultiplyAdd,Negate,NextDouble,Norm,OpenRead,OpenWrite,OrderBy,
    OrderByDescending,Parse,Peek,Pop,Pow,ProductOfSums,Push,Range,
    ReadAllLines,ReadAllLinesAsync,ReadAllText,ReadAllTextAsync,
    ReadAsync,ReadDouble,ReadInt32,ReadDoubleLittleEndian,Remove,
    RemoveAll,RemoveAt,Rent,Replace,Return,Reverse,RootN,Round,Select,
    SelectMany,SequenceEqual,Serialize,SerializeAsync,Sign,Sin,Sinh,
    Sigmoid,Skip,Slice,Sort,Sqrt,StoreUnsafe,Subtract,Sum,SumOfSquares,
    SuppressFinalize,Take,TakeLast,Tan,Tanh,ThenBy,ThrowIfNull,
    ThrowIfNegative,ThrowIfNegativeOrZero,ThrowIfGreaterThan,
    ToArray,ToDictionary,ToHashSet,ToList,ToLookup,ToString,Truncate,
    TryAdd,TryDequeue,TryFormat,TryGetValue,TryParse,UnionWith,
    WaitAll,WhenAll,WhenAny,Where,Write,WriteAllLines,WriteAllText,
    WriteAllTextAsync,WriteAsync,WriteUnaligned,
    WriteDoubleLittleEndian,WriteLine,Zip,CreateDirectory,
    GetExtension,GetField,GetFileNameWithoutExtension,GetFullPath,
    StartActivity,StartNew,Stop,Trim,TrimStart,TrimEnd,Compare,
    OfType,LastOrDefault,Single,SingleOrDefault,ElementAt,Distinct,
    OrderByDescending,Aggregate,IsNullOrEmpty,IsNullOrWhiteSpace,
    Join,Substring,StartsWith,EndsWith,ToLower,ToUpper,LoadAsync,
    Process,Compute,Run,SuppressFinalize,Step,Solve,Work,
    Exists,GetLength,FromDateTime},
  morecomment=[l]{//},
  morecomment=[s]{/*}{*/},
  morestring=[b]",
  morestring=[b]',
}
\lstdefinestyle{cs}{
  language=CSharp,
  basicstyle=\ttfamily\bfseries\fontsize{6.60pt}{9.4pt}\selectfont,
  keywordstyle=[1]{\bfseries\color{clKw}},
  keywordstyle=[2]{\color{clTy}},
  keywordstyle=[3]{\color{clFn}},
  stringstyle=\color{clSt},
  commentstyle=\itshape\color{clCm},
  backgroundcolor=\color{bgCode},
  frame=none,
  xleftmargin=3pt, xrightmargin=2pt,
  aboveskip=2pt, belowskip=1pt,
  breaklines=true, keepspaces=true,
  showstringspaces=false, tabsize=2,
  literate=
    {=>}{{\textbf{\color{clOp}=>}}}{2}
    {->}{{\textbf{\color{clOp}->}}}{2}
    {??}{{\textbf{\color{clOp}??}}}{2}
    {?.}{{\textbf{\color{clOp}?.}}}{2}
    {>=}{{\textbf{\color{clOp}>=}}}{2}
    {<=}{{\textbf{\color{clOp}<=}}}{2}
    {==}{{\textbf{\color{clOp}==}}}{2}
    {!=}{{\textbf{\color{clOp}!=}}}{2},
}

\usepackage{tcolorbox}
\tcbuselibrary{skins,breakable}
\newenvironment{card}[3][acBlue]{%
  \begin{tcolorbox}[
    enhanced, colback=#2, colframe=bdCard,
    coltitle=black, colbacktitle=#2!70!bgHdr,
    fonttitle=\bfseries\fontsize{6.16pt}{7.48pt}\selectfont\MakeUppercase,
    title={#3},
    left=3pt, right=2pt, top=1.4pt, bottom=1.8pt,
    toptitle=1pt, bottomtitle=1pt,
    boxrule=0.4pt,
    borderline west={3pt}{0pt}{#1},
    arc=2pt, before skip=0pt, after skip=1.6mm,
    breakable,
  ]}{\end{tcolorbox}}

\usepackage{multicol}
\setlength{\columnsep}{2.6mm}

\newcommand{\sectionbanner}[1]{%
  \vspace{0.4mm}%
  {\setlength{\fboxsep}{1.5pt}%
   \colorbox{secBg}{\color{white}\bfseries\fontsize{6.6pt}{7.9pt}\selectfont
   \MakeUppercase{\hspace{0.5mm}#1\hspace{0.5mm}}}%
  }\vspace{1.3mm}\par\nobreak}

\newcommand{\sub}[1]{%
  \par\vspace{1.5pt}%
  {\fontsize{5.7pt}{7pt}\selectfont\color{gray}\MakeUppercase{#1}}\par
  \vspace{0.5pt}}

\newcommand{\code}[1]{{\ttfamily\bfseries #1}}
\newcommand{\trap}[1]{{\fontsize{5.6pt}{6.8pt}\selectfont\color{trap}\bfseries TRAP:\,#1}}
\newcommand{\note}[1]{{\fontsize{5.6pt}{6.8pt}\selectfont\itshape #1}}

\usepackage{booktabs,array,colortbl,enumitem}
\setlist[itemize]{leftmargin=3.5mm,itemsep=0.5pt,topsep=1pt,
                  parsep=0pt,partopsep=0pt}

% =====================================================================
\begin{document}
% =====================================================================

{\centering
{\fontsize{13pt}{15.6pt}\selectfont\bfseries
  P. Areias simple C\#}\par\vspace{0.6mm}
{\fontsize{6.4pt}{7.6pt}\selectfont\color{gray}
  \MakeUppercase{Numerical \& Scientific Computing Edition}}\par\vspace{0.8mm}
{\color{secBg}\rule{\linewidth}{1pt}}\par\vspace{1.5mm}}

% =====================================================================
\sectionbanner{Language Fundamentals}
\begin{multicols}{2}

\begin{card}[acBlue]{bgBlue}{Program Structure \& Namespaces}
\begin{lstlisting}[style=cs]
// File-scoped namespace is fine; keep the entry point explicit.
using System;
using System.Numerics;
using static System.Math;       // Sqrt(), Sin(), PI

namespace MyApp.Numerics;

using DoubleVector = double[];  // local type alias

public static class Program
{
    public static int Main(string[] args)
    {
        Console.WriteLine($"args = {args.Length}");
        double[] x = [1.0, 2.0, 3.0];
        Console.WriteLine($"norm = {Norm(x):G17}");
        return 0;
    }

    private static double Norm(ReadOnlySpan<double> x)
    {
        double s = 0.0;
        for (int i = 0; i < x.Length; i++)
            s += x[i] * x[i];
        return Sqrt(s);
    }
}

internal sealed class InternalHelper { }
\end{lstlisting}
\sub{Build, Run \& Project Switches}
\begin{lstlisting}[style=cs]
// $ dotnet new console -n MyApp
// $ cd MyApp && dotnet run -c Release
// $ dotnet publish -c Release -r linux-x64 --self-contained

// .csproj essentials for strict numerical code
// <TargetFramework>net10.0</TargetFramework>
// <LangVersion>14.0</LangVersion>   // avoid latest for reproducibility
// <Nullable>enable</Nullable>
// <ImplicitUsings>enable</ImplicitUsings>
// <CheckForOverflowUnderflow>true</CheckForOverflowUnderflow>
// <AllowUnsafeBlocks>true</AllowUnsafeBlocks>  // only when needed
\end{lstlisting}
\end{card}

\begin{card}[acTeal]{bgTeal}{C\# 14 Essentials}
\begin{lstlisting}[style=cs]
// Null-conditional assignment: RHS evaluated only if c != null
customer?.Order = GetCurrentOrder();
customer?.Count += 1;      // compound assignment is allowed
// customer?.Count++;      // not allowed

// nameof with an unbound generic type
string name = nameof(Dictionary<,>);    // "Dictionary"

// Extension members add instance-like members from static classes
// Use sparingly: prefer obvious APIs for numerical kernels.

// Field-backed properties: 'field' is contextual
public string Name {
  get;
  set => field = value ?? throw new ArgumentNullException(nameof(value));
}

// Partial constructors/events and user-defined compound assignment
// are mainly library/source-generator features; keep public APIs simple.
\end{lstlisting}
\note{C\# 14 is tied to the .NET 10 SDK in normal SDK-style projects.
Prefer the target framework's default language version or pin
\code{14.0}; avoid \code{latest} in reproducible builds.}
\end{card}

\begin{card}[acNavy]{bgNavy}{Built-in Types Reference}
\fontsize{6.0pt}{7.4pt}\selectfont
\begin{tabular}{@{}p{16mm}p{32mm}@{}}
\rowcolor{bgHdr}\textbf{Type} & \textbf{Notes}\\
\code{bool}     & \code{true} / \code{false}\\
\code{byte}, \code{sbyte} & 8-bit unsigned / signed\\
\code{short}, \code{ushort} & 16-bit\\
\code{int}, \code{uint} & 32-bit; default for int literals\\
\code{long}, \code{ulong} & 64-bit\\
\code{Int128}, \code{UInt128} & 128-bit\\
\code{nint}, \code{nuint} & native size (32 on x86, 64 on x64)\\
\code{Half}     & 16-bit IEEE 754; \code{System.Half}\\
\code{float}    & 32-bit IEEE 754; suffix \code{f}\\
\code{double}   & 64-bit IEEE 754 -- \textbf{default for FEM}\\
\code{decimal}  & 128-bit base-10; for money, \emph{not} science\\
\code{Complex}  & \code{System.Numerics.Complex}\\
\code{char}     & UTF-16 code unit; \code{string}: immutable seq.\\
\end{tabular}
\sub{IEEE 754 essentials}
\begin{lstlisting}[style=cs]
double inf = double.PositiveInfinity;
double nan = double.NaN;
double eps = double.Epsilon;          // smallest +denormal
const double Macheps = 2.220446049250313e-16; // nextafter(1)-1
const double UnitRoundoff = 1.1102230246251565e-16; // eps/2

bool ok = double.IsFinite(x);   // not NaN, not Infinity
bool sub = double.IsSubnormal(x);
\end{lstlisting}
\trap{\code{x == double.NaN} is always \code{false}; use \code{double.IsNaN(x)}.}
\trap{\code{double.NaN.CompareTo(1.0)} returns $-1$ (NaN sorts as least);
  but \code{NaN < 1.0} is \code{false}. Sort and compare disagree.}
\end{card}

\begin{card}[acOrange]{bgOrange}{System.Math, double \& Numeric Constants}
\fontsize{6.05pt}{7.5pt}\selectfont
\begin{tabular}{@{}p{24mm}p{24mm}@{}}
\rowcolor{bgHdr}\textbf{Method} & \textbf{Notes}\\
\code{Sin Cos Tan} & radians\\
\code{Asin Acos Atan} & inverses\\
\code{Atan2(y, x)} & full quadrant\\
\code{Sinh Cosh Tanh} & hyperbolic\\
\code{Sqrt Cbrt}  & roots\\
\code{Pow(a, b)}  & $a^b$\\
\code{Exp Log Log2 Log10} & exp \& logs\\
\code{Abs Sign}   & magnitude / $\pm 1$\\
\code{Max Min Clamp} & $v \in [lo, hi]$\\
\code{Ceiling Floor Truncate} & rounding\\
\code{Round(x, n, mode)} & banker's by default\\
\code{FusedMultiplyAdd(a,b,c)} & $a\!\cdot\! b + c$ one rounding\\
\code{double.Hypot(x,y)} & $\sqrt{x^2+y^2}$ overflow-safe\\
\code{ScaleB(x, n)} & $x \cdot 2^n$ exact\\
\code{ILogB(x)}     & integer $\lfloor \log_2 |x| \rfloor$\\
\code{BitDecrement(x)} & next double toward $-\infty$\\
\code{BitIncrement(x)} & next double toward $+\infty$\\
\end{tabular}
\sub{Constants \& precision notes}
\begin{lstlisting}[style=cs]
double pi = Math.PI;     // 3.14159265358979...
double tau = Math.Tau;   // 2 * PI
\end{lstlisting}
\note{\code{Math} is double; \code{MathF} is float. Static generic math
members such as \code{double.Hypot} and \code{T.Hypot} complement
\code{Math}. \code{FusedMultiplyAdd} gives one rounding instead of two;
it may be faster when hardware and the JIT select fused instructions.}
\end{card}

\begin{card}[acRed]{bgRed}{Operators \& Pattern Matching}
\fontsize{6.55pt}{8.6pt}\selectfont
\begin{tabular}{@{}p{18mm}@{\;\;}p{27mm}@{}}
\rowcolor{bgHdr}\textbf{Expr} & \textbf{Meaning}\\
{\ttfamily\color{clOp}\bfseries !\;||\;\&\&} & not / or / and (short-circuit)\\
{\ttfamily\color{clOp}\bfseries ??\;??=}    & null-coalescing / assign\\
{\ttfamily\color{clOp}\bfseries ?.}\;{\ttfamily\bfseries ?[]} & null-conditional access\\
{\ttfamily\color{clOp}\bfseries ?.X = v}    & null-cond.\ assign\\
{\ttfamily\color{clKw}\bfseries is \color{clTy}T} {\ttfamily\bfseries v} & type test + bind\\
{\ttfamily\color{clKw}\bfseries is not null} & non-null check\\
{\ttfamily\color{clKw}\bfseries is > 0}      & relational pattern\\
{\ttfamily\color{clKw}\bfseries is (>0 and <}{\ttfamily\bfseries N)} & combined patterns\\
{\ttfamily\color{clKw}\bfseries is [}{\ttfamily\bfseries x, ..]} & list/slice pattern\\
{\ttfamily\bfseries x \color{clKw} as \color{clTy}T} & safe cast; null on fail\\
{\ttfamily\color{clKw}\bfseries sizeof(T)}   & bytes (unmanaged \code{T})\\
{\ttfamily\color{clKw}\bfseries nameof(D<,>)} & string name; unbound\\
{\ttfamily\color{clKw}\bfseries checked\{\}} & overflow throws\\
\end{tabular}
\begin{lstlisting}[style=cs]
// Switch expression: exhaustive, returns value
double abs = x switch {
  double.NaN => throw new ArithmeticException(), // test first
  >= 0       => x,
  < 0        => -x,
  _          => throw new ArithmeticException()
};
// Property pattern + tuple pattern
static int Quad((double X, double Y) p) => p switch {
  { X: > 0, Y: > 0 } => 1,   { X: < 0, Y: > 0 } => 2,
  { X: < 0, Y: < 0 } => 3,   { X: > 0, Y: < 0 } => 4,
  _                  => 0
};
\end{lstlisting}
\end{card}

\begin{card}[acGreen]{bgGreen}{Control Flow}
\begin{lstlisting}[style=cs]
if (x > 0) { } else if (x < 0) { } else { }
for (int i = 0; i < n; i++) { break; continue; }
foreach (var v in seq) { }       // cannot mutate seq
while (cond) { }
do { } while (cond);

// Switch statement: explicit break required
switch (kind) {
  case 1:          A(); break;
  case 2: case 3:  B(); break;
  case int v when v > 10:   C(v); break;
  default:         D(); break;
}

// Static local: no captures, no allocation
double Solve(double[] x) {
  static double Inv(double v) =>
      v == 0 ? throw new DivideByZeroException() : 1/v;
  return Inv(x[0]);
}
\end{lstlisting}
\end{card}

\begin{card}[acRed]{bgRed}{Exception Handling}
\begin{lstlisting}[style=cs]
try {
  Solve(K, rhs, x);
}
catch (FileNotFoundException ex) {
  Log($"missing: {ex.FileName}");
}
catch (ArithmeticException ex) when (ex.Message.Contains("nan")) {
  Recover();                 // exception filter
}
catch (Exception) {
  throw;                     // re-throw; keeps stack trace
                             // do NOT write 'throw ex;'
}
finally { /* always runs; release locks here */ }

// Guard helpers (CallerArgumentExpression makes good msgs)
ArgumentNullException.ThrowIfNull(coords);
ArgumentOutOfRangeException.ThrowIfNegative(n);
ArgumentOutOfRangeException.ThrowIfGreaterThan(i, max);

// Throw expression
var v = map.ContainsKey(k) ? map[k]
      : throw new KeyNotFoundException(nameof(k));
\end{lstlisting}
\note{Exceptions are for exceptional paths. Validate at the public API
boundary, avoid exception-driven control flow in kernels, and keep hot loops
free of guards once preconditions have been checked.}
\end{card}

\begin{card}[acPurple]{bgPurple}{Arrays, Ranges \& Collection Expr.}
\begin{lstlisting}[style=cs]
double[]   a   = new double[n];      // zero-initialised
double[,]  mat = new double[r, c];   // 2-D rectangular
double[][] jag = new double[r][];    // jagged rows

// Index ^ and range ..
double   last = a[^1];        // a[a.Length - 1]
double[] mid  = a[2..^2];     // ALLOCATES (new array)
Span<double> view = a.AsSpan(2, n - 4); // does NOT allocate

// Bulk operations -- prefer over hand loops
Array.Copy(src, sOff, dst, dOff, count);
Array.Clear(a);                  // zeroes every element
Array.Sort(keys, values);        // co-sort parallel arrays

// Collection expressions
int[]     p = [1, 2, 3];
List<int> q = [..p, 4, 5];       // spread operator
double[]  r = [..a, ..b];        // concatenate

// Named tuples
(double Lo, double Hi) MinMax(ReadOnlySpan<double> x) {
  double lo = x[0], hi = x[0];
  for (int i = 1; i < x.Length; i++) {
    if (x[i] < lo) lo = x[i];
    if (x[i] > hi) hi = x[i];
  }
  return (lo, hi);
}
var (lo, hi) = MinMax(data);    // deconstruct
\end{lstlisting}
\trap{\code{a[2..8]} on an array allocates a new array. Inside loops,
always use \code{a.AsSpan(2, 6)}.}
\end{card}

\end{multicols}

% =====================================================================
\sectionbanner{Types \& Object-Oriented Programming}
\begin{multicols}{2}

\begin{card}[acBlue]{bgBlue}{Class --- Modern Numerical Style}
\begin{lstlisting}[style=cs]
public sealed class Mesh(double[] coords) : IEquatable<Mesh>
{
  private readonly double[] _coords =
      coords ?? throw new ArgumentNullException(nameof(coords));

  public required string Name { get; init; }  // set by initializer
  public required int    Dim  { get; init; }

  public IReadOnlyList<double> Coords => _coords;

  // Field-backed property: 'field' is compiler-generated storage.
  public int NodeCount {
    get;
    set => field = value < 0
      ? throw new ArgumentOutOfRangeException(nameof(value))
      : value;
  } = coords.Length / 3;

  public bool Equals(Mesh? other) =>
      other is not null && Name == other.Name && Dim == other.Dim;

  public override bool Equals(object? obj) => obj is Mesh m && Equals(m);
  public override int GetHashCode() => HashCode.Combine(Name, Dim);

  public static bool operator==(Mesh? a, Mesh? b) =>
      a is null ? b is null : a.Equals(b);
  public static bool operator!=(Mesh? a, Mesh? b) => !(a == b);
}

var mesh = new Mesh(coords) { Name = "beam", Dim = 3 };
\end{lstlisting}
\note{Primary constructors place parameters in scope for instance initializers
and members. They do not imply ownership: copy or validate mutable inputs when
object invariants require it. Equal objects must hash equal, and hash values
must remain stable while used as dictionary keys.}
\end{card}

\begin{card}[acTeal]{bgTeal}{Records \& Record Structs (Vec3)}
\begin{lstlisting}[style=cs]
// record class: reference type, value-based equality
// Compiler synthesises Equals, GetHashCode, ToString,
// operator==, copy ctor, Deconstruct, with-expressions.
public record Point(double X, double Y) {
  public double Norm => Math.Sqrt(X*X + Y*Y);
}
Point a = new(1, 2);
Point b = a with { Y = 5 };       // non-destructive copy
var (x, y) = a;                   // deconstruct

// readonly record struct: value type; no separate heap object
// unless boxed/captured. Good for tiny vectors in inner loops
public readonly record struct Vec3(double X, double Y, double Z)
{
  public double Norm  => Math.Sqrt(X*X + Y*Y + Z*Z);
  public Vec3   Unit  => this / Norm;

  // Operator overloading: must be 'static'
  // BOTH scalar orderings must be defined
  public static Vec3 operator+(Vec3 a, Vec3 b) =>
      new(a.X+b.X, a.Y+b.Y, a.Z+b.Z);
  public static Vec3 operator-(Vec3 a)             // unary
      => new(-a.X, -a.Y, -a.Z);
  public static Vec3 operator*(double s, Vec3 v) =>
      new(s*v.X, s*v.Y, s*v.Z);
  public static Vec3 operator*(Vec3 v, double s) => s * v;
  public static Vec3 operator/(Vec3 v, double s) =>
      new(v.X/s, v.Y/s, v.Z/s);

  public static double Dot(Vec3 a, Vec3 b) =>
      a.X*b.X + a.Y*b.Y + a.Z*b.Z;
  public static Vec3 Cross(Vec3 a, Vec3 b) => new(
      a.Y*b.Z - a.Z*b.Y, a.Z*b.X - a.X*b.Z,
      a.X*b.Y - a.Y*b.X);
}
\end{lstlisting}
\note{Mark structs \code{readonly} when possible: this avoids defensive
copies on member access through \code{in} parameters. Value types are not
``always stack allocated''; they may live inline in arrays or objects.}
\end{card}

\begin{card}[acPurple]{bgPurple}{Generics, INumber\textlangle T\textrangle\ \& Constraints}
\begin{lstlisting}[style=cs]
// where T : class | struct | unmanaged | new()
//       | IComparable<T> | INumber<T> | IFloatingPointIeee754<T>
//       | allows ref struct  -- accepts Span<T> etc.

// INumberBase<T> static surface (numeric primitives expose):
//   T.Zero, T.One, T.AdditiveIdentity, T.MultiplicativeIdentity
//   T.Abs(x), T.Min(a,b), T.Max(a,b), T.Sign(x)
//   T.IsFinite(x), T.IsNaN(x), T.IsZero(x), T.IsNegative(x)
//   T.CreateChecked<U>(x)    -- throws OverflowException
//   T.CreateSaturating<U>(x) -- clamps to MinValue/MaxValue
//   T.CreateTruncating<U>(x) -- bit-truncates / wraps

// IFloatingPointIeee754<T> adds:
//   T.Pi, T.E, T.Tau, T.Epsilon, T.NaN, T.PositiveInfinity
//   T.Sqrt(x), T.Cbrt(x), T.Sin(x), T.Cos(x), T.Atan2(y,x)
//   T.Exp(x), T.Log(x), T.Pow(x,y), T.Hypot(x,y)
//   T.FusedMultiplyAdd(a,b,c)

public static T Norm<T>(ReadOnlySpan<T> x)
    where T : IFloatingPointIeee754<T>
{
  T s = T.Zero;
  for (int i = 0; i < x.Length; i++)
    s = T.FusedMultiplyAdd(x[i], x[i], s);  // fused per term
  return T.Sqrt(s);
}

// Generic class with ref-returning indexer
public sealed class DenseMatrix<T> where T : struct, INumber<T>
{
  private readonly T[] _d;
  public int Rows { get; }
  public int Cols { get; }
  public DenseMatrix(int r, int c)
      { Rows = r; Cols = c; _d = new T[r * c]; }

  // ref T indexer: K[i,j] += dK is a real in-place update
  public ref T this[int i, int j] => ref _d[i * Cols + j];
  public Span<T> Row(int i) => _d.AsSpan(i * Cols, Cols);
  public Span<T> AsFlat() => _d.AsSpan();
}
\end{lstlisting}
\end{card}

\begin{card}[acOrange]{bgOrange}{Extension Members \& Lambdas}
\begin{lstlisting}[style=cs]
// CLASSIC extension methods (still works everywhere)
public static class VectorExt {
  public static double L2(this double[] v) {
    double s = 0;
    for (int i = 0; i < v.Length; i++) s += v[i]*v[i];
    return Math.Sqrt(s);
  }
}

// EXTENSION MEMBERS: properties, operators, statics
public static class SpanExt
{
  // Instance extension on a generic receiver named 'x'
  extension<T>(ReadOnlySpan<T> x) where T : INumber<T>
  {
    public bool IsEmpty => x.Length == 0;       // property
    public T Sum() {
      T s = T.Zero;
      for (int i = 0; i < x.Length; i++) s += x[i];
      return s;
    }
  }
  // Static extension: adds a static factory to Span<T>
  extension<T>(Span<T>) {
    public static Span<T> Empty => default;
  }
  // Static extension on a numeric type
  extension(double) {
    public static double TwoPi => 2 * Math.PI;
  }
}
\end{lstlisting}
\sub{Lambdas, Func, Action}
\begin{lstlisting}[style=cs]
Func<double, double> f      = x => x * x;
Action<string>       print  = Console.WriteLine;
Predicate<double>    small  = x => Math.Abs(x) < 1e-12;
var inc = double (double x) => x + 1.0;       // typed

// Closures capture VARIABLES by reference, not values
double tol = 1e-10;
Predicate<double> tiny = x => Math.Abs(x) < tol;
tol = 1e-6; // 'tiny' now uses 1e-6
\end{lstlisting}
\end{card}

\begin{card}[acAmber]{bgAmber}{Enums \& Flags}
\begin{lstlisting}[style=cs]
public enum BcKind   { Dirichlet, Neumann, Robin, Periodic }
public enum Solver : byte { CG = 1, GMRES = 2, DirectLU = 3 }

BcKind k = Enum.Parse<BcKind>("Neumann");
bool ok  = Enum.IsDefined<BcKind>(k);
foreach (BcKind v in Enum.GetValues<BcKind>()) { }

[Flags] public enum Dof
{
  None = 0,
  Ux = 1<<0, Uy = 1<<1, Uz = 1<<2,
  Rx = 1<<3, Ry = 1<<4, Rz = 1<<5,
  AllT = Ux | Uy | Uz,
  All  = AllT | Rx | Ry | Rz
}
Dof  pin    = Dof.Ux | Dof.Uy | Dof.Uz;
bool fixedX = (pin & Dof.Ux) != 0;       // hot-path style
pin |= Dof.Rz;                            // add a flag
pin &= ~Dof.Uy;                           // remove a flag
\end{lstlisting}
\note{In hot paths use bitwise tests, not \code{HasFlag} -- the bitwise
form is also clearer at a glance.}
\end{card}

\begin{card}[acRed]{bgRed}{Properties, Indexers \& Events}
\begin{lstlisting}[style=cs]
public sealed class Sensor
{
  // Init-only auto-property
  public string Id { get; init; } = "";

  // Computed -- no backing field
  public double Magnitude => Math.Sqrt(X*X + Y*Y);

  // Field-backed validating property
  public double X {
    get;
    set => field = double.IsFinite(value)
      ? value : throw new ArgumentException("not finite");
  }
  public double Y { get; set; }

  // Indexer
  public double this[int axis] => axis switch {
    0 => X, 1 => Y,
    _ => throw new IndexOutOfRangeException()
  };

  // Event
  public event EventHandler<double>? Threshold;
  public void Update(double v) {
    X = v;
    if (Math.Abs(v) > 1e6) Threshold?.Invoke(this, v);
  }
}
\end{lstlisting}
\end{card}

\begin{card}[acGreen]{bgGreen}{Nullable Reference Types}
\begin{lstlisting}[style=cs]
#nullable enable

string  a = "ok";          // non-null by contract
string? b = MaybeText();    // may be null

if (b is { Length: > 0 })
  Console.WriteLine(b.Length);   // flow analysis proves non-null

// Prefer explicit nullable return values over side-effect outputs
Node? node = Lookup(id);
if (node is not null)
  Use(node);

// Null-forgiving operator: use only when you know better than flow analysis
Node root = possiblyNullRoot!;
\end{lstlisting}
\note{Nullable reference types are compile-time analysis, not a runtime
non-null guarantee. In new SDK projects \code{<Nullable>enable</Nullable>}
is normally on; keep it on for public numerical libraries.}
\end{card}

\begin{card}[acNavy]{bgNavy}{IDisposable, IAsyncDisposable, Lock}
\begin{lstlisting}[style=cs]
public sealed class NativeBuffer : IDisposable
{
  private unsafe void* _p;
  private readonly int _count;

  public NativeBuffer(int count) {
    if (count < 0) throw new ArgumentOutOfRangeException(nameof(count));
    _count = count;
    nuint bytes = checked((nuint)count * (nuint)sizeof(double));
    unsafe { _p = NativeMemory.AlignedAlloc(bytes, alignment: 64); }
  }

  public unsafe Span<double> AsSpan() => new((double*)_p, _count);

  public unsafe void Dispose() {
    if (_p is null) return;
    NativeMemory.AlignedFree(_p);
    _p = null;
    GC.SuppressFinalize(this);
  }
}
using var buf = new NativeBuffer(8192);

// Async cleanup
public sealed class AsyncWriter : IAsyncDisposable {
  public async ValueTask DisposeAsync() {
    await _stream.FlushAsync();
    _stream.Dispose();
  }
}
await using var w = new AsyncWriter();

// Modern lock primitive -- prefer over lock(object)
private readonly Lock _gate = new();
using (_gate.EnterScope()) { /* critical section */ }
\end{lstlisting}
\note{If you write a finalizer, also call \code{GC.SuppressFinalize(this)}
in \code{Dispose} to avoid the queue overhead when the user disposes
deterministically.}
\end{card}

\end{multicols}

% =====================================================================
\sectionbanner{Collections}
\begin{multicols}{2}

\begin{card}[acBlue]{bgBlue}{List, Dictionary, HashSet}
\begin{lstlisting}[style=cs]
// PRE-SIZE when count is known; resize copies are silent killers
var nodes = new List<int>(capacity: nNodes);
List<int> ids = [1, 2, 3];          // collection expr

nodes.Add(n);  nodes.AddRange(batch);
nodes.RemoveAll(x => x < 0);        // single-pass filter
nodes.EnsureCapacity(2 * nNodes);
Span<int> rawNodes = CollectionsMarshal.AsSpan(nodes); // no List resize while used

// Dictionary -- O(1) keyed lookup
var w = new Dictionary<int, double>(capacity: 256);
w[id] = value;                       // add or overwrite
w.TryAdd(id, value);                 // false on duplicate
if (w.ContainsKey(id)) { double v = w[id]; } // two lookups, clear API

// CollectionsMarshal.AsSpan(List<T>) is fast but invalidated by resize.
// For dictionaries, prefer ordinary APIs unless profiling demands otherwise.

// HashSet for repeated membership tests
var seen = new HashSet<int>(capacity: 1024);
bool isNew = seen.Add(id);           // false if duplicate
seen.UnionWith(other); seen.IntersectWith(other);

// Sorted (Red-Black tree; O(log n))
var sd = new SortedDictionary<int, string>();

// PriorityQueue: min-heap on TPriority
var pq = new PriorityQueue<string, double>();
pq.Enqueue("task", priority: 3.0);
string item = pq.Dequeue();             // removes minimum priority item
\end{lstlisting}
\trap{Views returned by \code{CollectionsMarshal.AsSpan(List<T>)} are invalidated
by any operation that may resize the list. Use such views only in tight,
straight-line code where collection capacity is stable.}
\end{card}

\begin{card}[acTeal]{bgTeal}{Async Streams \& Channels}
\begin{lstlisting}[style=cs]
using System.Threading.Channels;

// Producer/consumer pipeline; bounded back-pressure
var ch = Channel.CreateBounded<double>(capacity: 1024);

_ = Task.Run(async () => {                  // producer
  for (int i = 0; i < n; i++)
    await ch.Writer.WriteAsync(i * h);
  ch.Writer.Complete();
});

await foreach (double v in ch.Reader.ReadAllAsync())
  Process(v);                               // consumer

// Async iterator method
public static async IAsyncEnumerable<string>
    StreamLinesAsync(string path)
{
  using var sr = new StreamReader(path);
  string? line;
  while ((line = await sr.ReadLineAsync()) is not null)
    yield return line;
}

// Lock-free MPMC queue; consumer APIs use the BCL Try-pattern.
// Hide that detail behind a small adapter if you dislike out parameters.
var q = new ConcurrentQueue<int>();
q.Enqueue(42);
\end{lstlisting}
\end{card}


\begin{card}[acPurple]{bgPurple}{LINQ: Useful, But Know the Cost}
\begin{lstlisting}[style=cs]
// Declarative and clear; often fine outside kernels
var top = nodes.Where(i => mass[i] > 0.0)
               .OrderBy(i => x[i])
               .Take(10)
               .ToArray();

// Hot-loop version: explicit, predictable, no iterator/delegate chain
int m = 0;
for (int k = 0; k < nodes.Count; k++) {
  int i = nodes[k];
  if (mass[i] > 0.0) scratch[m++] = i;
}
Array.Sort(scratch, 0, m, Comparer<int>.Create((i, j) => x[i].CompareTo(x[j])));
\end{lstlisting}
\note{LINQ allocates iterator/delegate objects in many cases and can hide
multiple passes. Use it freely for setup code; prefer explicit loops in
low-level numerical kernels and parsers.}
\end{card}
\end{multicols}

% =====================================================================
\sectionbanner{Spans, Memory \& High-Performance Code}
\begin{multicols}{2}

\begin{card}[acAmber]{bgAmber}{Span, Memory \& stackalloc}
\begin{lstlisting}[style=cs]
// Span<T>: stack-only ref-struct; zero-copy view; sync only
// CANNOT: be a class field, be boxed, cross await.
double[] arr = new double[1024];
Span<double> row = arr.AsSpan(off, len);   // no allocation
Span<double> tmp = stackalloc double[16];  // stack buffer
tmp.Clear();

// Implicit conversions: T[] -> Span<T> at call sites
static double Dot(ReadOnlySpan<double> a, ReadOnlySpan<double> b)
{
  if (a.Length != b.Length) throw new ArgumentException();
  double s = 0;
  for (int i = 0; i < a.Length; i++) s += a[i] * b[i];
  return s;
}
double d = Dot(arr, arr);   // direct, no .AsSpan() needed

// scoped: prevents argument from escaping caller
static int CountPos(scoped ReadOnlySpan<double> x) {
  int c = 0;
  for (int i = 0; i < x.Length; i++) if (x[i] > 0) c++;
  return c;
}

// Memory<T>: heap-backed; survives async; can be a field
Memory<double> mem = new double[n];
await ProcessAsync(mem);              // Span<T> CANNOT do this
Span<double> s = mem.Span;            // extract inside sync code

// params Span<T>: variadic without array allocation
static double Max(params ReadOnlySpan<double> values) {
  double m = double.NegativeInfinity;
  foreach (double v in values) if (v > m) m = v;
  return m;
}
\end{lstlisting}
\trap{\code{a[2..8]} on a \code{T[]} allocates a new array.
Inside loops, ALWAYS use \code{a.AsSpan(2, 6)}.}
\end{card}

\begin{card}[acRed]{bgRed}{ref Returns, ref Locals \& ref struct}
\begin{lstlisting}[style=cs]
// ref return + ref local: writes flow back to the source array
static ref double At(double[] a, int i) => ref a[i];
ref double cell = ref At(arr, 10);
cell *= 2;                          // updates arr[10]

// ref readonly return: alias without copy, no mutation
static ref readonly Vec3 First(ReadOnlySpan<Vec3> x)
    => ref x[0];

// Prefer value-returning APIs over ref-parameter APIs.
static Vec3 Negated(Vec3 v) => new(-v.X, -v.Y, -v.Z);

static double ParseStrict(string s) =>
    double.Parse(s, CultureInfo.InvariantCulture);

// BCL Try* APIs use out; consume them when needed,
// but avoid designing your own numerical APIs around out parameters.

// ref struct (Span<T>, ReadOnlySpan<T>, custom): stack-only
public ref struct StackBuf {
  private Span<double> _data;
  public StackBuf(Span<double> d) => _data = d;
  public ref double this[int i] => ref _data[i];
}
\end{lstlisting}
\trap{Avoid \code{ref} parameters in public numerical APIs unless a measured
interop or ABI constraint forces them. \code{ref} returns and short-lived
\code{ref} locals are useful for indexed storage views; do not let such
aliases outlive the storage operation that produced them.}
\end{card}

\begin{card}[acOrange]{bgOrange}{ArrayPool, NativeMemory, GC Control}
\begin{lstlisting}[style=cs]
// Rent/return: avoids GC pressure for large transient buffers
var pool = ArrayPool<double>.Shared;
double[] buf = pool.Rent(n);   // may be larger than n!
try {
  Span<double> s = buf.AsSpan(0, n);
  s.Clear();                      // never assume zeroed
  Work(s);
}
finally { pool.Return(buf, clearArray: false); }

// Disposable wrapper for clean call-sites
public sealed class Scratch(int n) : IDisposable {
  private double[]? _b = ArrayPool<double>.Shared.Rent(n);
  public Span<double> Span => _b!.AsSpan(0, n);
  public void Dispose() {
    if (_b is null) return;
    ArrayPool<double>.Shared.Return(_b);
    _b = null;
  }
}
using (var s = new Scratch(4096)) Work(s.Span);

// Skip zero-init when all elements will be overwritten
double[] big = GC.AllocateUninitializedArray<double>(n);

// Pinned heap allocation -- never moves; safe for native libs
double[] pin = GC.AllocateArray<double>(n, pinned: true);

// Aligned native allocation (32 = AVX2; 64 = AVX-512)
unsafe {
  double* p = (double*)NativeMemory.AlignedAlloc(
      (nuint)(n * sizeof(double)), alignment: 64);
  try { /* use p[0..n-1] */ }
  finally { NativeMemory.AlignedFree(p); }
}
\end{lstlisting}
\note{Arrays $> 85$\,kB land on the Large Object Heap (LOH), which is
collected only on Gen 2. Use \code{ArrayPool} or pinned allocation for
these to avoid LOH fragmentation.}
\end{card}

\begin{card}[acPurple]{bgPurple}{MemoryMarshal, Pointers \& Binary I/O}
\begin{lstlisting}[style=cs]
// Reinterpret a byte buffer as doubles (no copy)
Span<byte>   raw = stackalloc byte[64];
Span<double> dbl = MemoryMarshal.Cast<byte, double>(raw);

// Clear or fill through Span<T> first; the JIT usually removes
// redundant bounds checks in simple counted loops.
dbl.Clear();
for (int i = 0; i < dbl.Length; i++)
  dbl[i] = 0.0;

// Pin for raw pointer access in unsafe blocks
unsafe {
  fixed (double* p = dbl) { /* call native(p, ...) */ }
}

// Endian-aware binary I/O
using System.Buffers.Binary;
Span<byte> bytes = stackalloc byte[8];
BinaryPrimitives.WriteDoubleLittleEndian(bytes, 3.14);
double back = BinaryPrimitives.ReadDoubleLittleEndian(bytes);

// Ref to first element; not pinned by itself
ref double headRef = ref MemoryMarshal.GetArrayDataReference(arr);
\end{lstlisting}
\sub{Hot-loop discipline}
\begin{itemize}
\item Capturing lambdas/delegates allocate or extend lifetimes; non-capturing delegates may be cached.
\item Boxing via \code{object} or non-generic interfaces.
\item Property access on large mutable structs (defensive copy).
\item \code{string.Split} on huge inputs; use \code{SearchValues<char>}.
\end{itemize}
\sub{Closure capture trap}
\begin{lstlisting}[style=cs]
// BUG: every task captures the SAME 'i'
for (int i = 0; i < n; i++)
  tasks[i] = Task.Run(() => Work(i));   // all see final i

// FIX: introduce a fresh local
for (int i = 0; i < n; i++) {
  int j = i;
  tasks[i] = Task.Run(() => Work(j));
}
\end{lstlisting}
\end{card}


\begin{card}[acBlue]{bgBlue}{Native Interop: Keep Boundaries Explicit}
\begin{lstlisting}[style=cs]
using System.Runtime.InteropServices;

internal static partial class NativeKernels
{
  [LibraryImport("nativekernels", EntryPoint = "scale_inplace")]
  internal static unsafe partial void ScaleInPlace(
      double* data, int length, double factor);
}

// Keep managed arrays pinned only for the shortest possible scope.
unsafe {
  fixed (double* p = values) {
    NativeKernels.ScaleInPlace(p, values.Length, factor);
  }
}
\end{lstlisting}
\note{For new interop code, source-generated \code{LibraryImport} is usually
preferable to reflection-based \code{DllImport}. Validate calling convention,
array layout, ownership, and lifetime explicitly.}
\end{card}
\end{multicols}

% =====================================================================
\sectionbanner{SIMD \& Vectorised Numerics}
\begin{multicols}{2}

\begin{card}[acTeal]{bgTeal}{TensorPrimitives --- Use This First}
\begin{lstlisting}[style=cs]
// System.Numerics.Tensors: hardware-accelerated vector
// operations over Span<T>; auto-uses available SIMD where supported.
// NuGet/package may be needed depending on TFM.
using System.Numerics.Tensors;

ReadOnlySpan<double> a = ..., b = ..., c = ...;
Span<double> y = ...;

// Element-wise: y = a + b   (NOT matmul!)
TensorPrimitives.Add(a, b, y);
TensorPrimitives.Subtract(a, b, y);
TensorPrimitives.Multiply(a, b, y);        // Hadamard product
TensorPrimitives.Divide(a, b, y);

// Scalar/broadcast and axpy-style operations
TensorPrimitives.Add(a, 1.0, y);           // y = a + 1
TensorPrimitives.Multiply(a, 2.0, y);      // y = 2*a
TensorPrimitives.MultiplyAdd(a, 2.0, b, y);// y = 2*a + b, not fused
TensorPrimitives.FusedMultiplyAdd(a, b, c, y); // y = fma(a,b,c)

// Reductions
var sum   = TensorPrimitives.Sum<double>(a);
var sumSq = TensorPrimitives.SumOfSquares<double>(a);
var norm2 = TensorPrimitives.Norm<double>(a);   // sqrt(sumSq)
var dot   = TensorPrimitives.Dot<double>(a, b);
var dist  = TensorPrimitives.Distance<double>(a, b);
var cos   = TensorPrimitives.CosineSimilarity<double>(a, b);
var mxV   = TensorPrimitives.Max<double>(a);
var mxI   = TensorPrimitives.IndexOfMax<double>(a);

// Math (vectorised where supported)
TensorPrimitives.Abs(a, y);   TensorPrimitives.Sqrt(a, y);
TensorPrimitives.Exp(a, y);   TensorPrimitives.Log(a, y);
TensorPrimitives.Sin(a, y);   TensorPrimitives.Cos(a, y);
TensorPrimitives.Sigmoid(a, y); TensorPrimitives.SoftMax(a, y);
\end{lstlisting}
\note{Reach for \code{TensorPrimitives} before hand-written SIMD. It
handles tails and has tuned implementations. \code{MultiplyAdd} is the
ordinary multiply-plus-add API; use \code{FusedMultiplyAdd} for one-rounding
FMA semantics when available for the element type.}
\end{card}

\begin{card}[acAmber]{bgAmber}{Portable SIMD: Vector\textlangle T\textrangle}
\begin{lstlisting}[style=cs]
using System.Numerics;
// Vector<T> picks the widest hardware width at runtime
// (Vector128 / Vector256 / Vector512) -- count via Vector<T>.Count

static double DotSimd(ReadOnlySpan<double> a,
                      ReadOnlySpan<double> b)
{
  if (!Vector.IsHardwareAccelerated || a.Length < Vector<double>.Count)
    return DotScalar(a, b);

  int w = Vector<double>.Count;
  Vector<double> acc = Vector<double>.Zero;
  int i = 0;
  for (; i <= a.Length - w; i += w) {
    var va = new Vector<double>(a.Slice(i, w));
    var vb = new Vector<double>(b.Slice(i, w));
    acc += va * vb;
  }
  double s = Vector.Sum(acc);
  for (; i < a.Length; i++) s += a[i] * b[i];   // tail
  return s;
}
\end{lstlisting}
\sub{Width-explicit (AVX-512 path)}
\begin{lstlisting}[style=cs]
using System.Runtime.Intrinsics;
// Vector256: 4 doubles; Vector512: 8 doubles (AVX-512)
if (Vector512.IsHardwareAccelerated) {
  Vector512<double> acc = Vector512<double>.Zero;
  int i = 0;
  for (; i <= a.Length - 8; i += 8) {
    var va = Vector512.Create<double>(a.Slice(i, 8));
    var vb = Vector512.Create<double>(b.Slice(i, 8));
    acc = Vector512.Add(acc, Vector512.Multiply(va, vb));
  }
  double s = Vector512.Sum(acc);
  for (; i < a.Length; i++) s += a[i] * b[i];  // required tail
}
\end{lstlisting}
\note{Data layout and memory bandwidth often dominate instruction choice.
Use \code{NativeMemory.AlignedAlloc} when exact alignment matters; pinned
managed arrays are non-moving but should not be treated as a portable
64-byte-alignment guarantee.}
\end{card}

\begin{card}[acGreen]{bgGreen}{Generic Math Kernels --- Numerical Quality}
\begin{lstlisting}[style=cs]
// AXPY: y = alpha*x + y
public static void Axpy<T>(T alpha, ReadOnlySpan<T> x, Span<T> y)
    where T : INumber<T>
{
  if (x.Length != y.Length) throw new ArgumentException();
  for (int i = 0; i < x.Length; i++)
    y[i] += alpha * x[i];
}

// Kahan compensated sum: good when cancellation is moderate
public static T SumKahan<T>(ReadOnlySpan<T> x)
    where T : IFloatingPointIeee754<T>
{
  T s = T.Zero, c = T.Zero;
  for (int i = 0; i < x.Length; i++) {
    T y = x[i] - c;
    T t = s + y;
    c = (t - s) - y;
    s = t;
  }
  return s;
}

// Neumaier variant: better when |x[i]| can exceed |sum|
public static T SumNeumaier<T>(ReadOnlySpan<T> x)
    where T : IFloatingPointIeee754<T>
{
  T s = T.Zero, c = T.Zero;
  for (int i = 0; i < x.Length; i++) {
    T t = s + x[i];
    c += T.Abs(s) >= T.Abs(x[i]) ? (s - t) + x[i] : (x[i] - t) + s;
    s = t;
  }
  return s + c;
}

// Pairwise sum: simple, reproducible for fixed partitioning
public static T SumPairwise<T>(ReadOnlySpan<T> x) where T : INumber<T>
{
  if (x.Length <= 16) { T s = T.Zero; foreach (T v in x) s += v; return s; }
  int m = x.Length / 2;
  return SumPairwise(x[..m]) + SumPairwise(x[m..]); // span slices: no array
}

public static T Hypot<T>(T x, T y)
    where T : IFloatingPointIeee754<T> => T.Hypot(x, y);
\end{lstlisting}
\trap{Generic math can inhibit some specialization and inlining decisions.
For small dense kernels over \code{double}, keep specialized overloads and
profile against the generic version.}
\end{card}

\end{multicols}

% =====================================================================
\sectionbanner{Concurrency, Async \& Parallel}
\begin{multicols}{2}

\begin{card}[acAmber]{bgAmber}{async / await, Task \& ValueTask}
\begin{lstlisting}[style=cs]
// async returns Task / Task<T> / ValueTask / ValueTask<T>
public async Task<double[]> LoadAsync(string path,
    CancellationToken ct = default)
{
  string[] lines = await File.ReadAllLinesAsync(path, ct);
  return Array.ConvertAll(lines,
      s => double.Parse(s, CultureInfo.InvariantCulture));
}

double[] data = await LoadAsync("vec.txt");

// Run several awaits concurrently, then await the bundle
Task<double[]> t1 = LoadAsync("a.txt");
Task<double[]> t2 = LoadAsync("b.txt");
double[][] both = await Task.WhenAll(t1, t2);

// Cancellation: pass token everywhere; check periodically
var cts = new CancellationTokenSource();
cts.CancelAfter(TimeSpan.FromSeconds(30));
ct.ThrowIfCancellationRequested();

// ValueTask: avoids heap alloc when synchronously cached
public ValueTask<double> CachedAsync(int key) {
  if (_cache.ContainsKey(key))
    return ValueTask.FromResult(_cache[key]);
  return new ValueTask<double>(LoadFromDiskAsync(key));
}
\end{lstlisting}
\trap{Do NOT \code{await} the same \code{ValueTask} twice. Convert
to \code{Task} via \code{.AsTask()} if you need to.}
\trap{\code{Task.Run} is not a numerical parallel loop primitive. Use
\code{Parallel.For/ForEach}, partitioned tasks, SIMD, or \code{TensorPrimitives}
for CPU kernels; use \code{Task.Run} for UI/offload boundaries or unavoidable
synchronous work.}
\end{card}

\begin{card}[acNavy]{bgNavy}{Parallel \& Thread Safety}
\begin{lstlisting}[style=cs]
// Data parallelism: each iteration MUST be independent
Parallel.For(0, n, i => results[i] = Compute(input[i]));

var opt = new ParallelOptions {
  MaxDegreeOfParallelism = Math.Max(1, Environment.ProcessorCount - 1),
  CancellationToken      = cts.Token
};

// Thread-safe reduction with thread-local accumulator
object sumLock = new();
double total = 0.0;
Parallel.For<double>(0, n,
  () => 0.0,
  (i, _, local) => local + values[i] * values[i],
  local => { lock (sumLock) total += local; });

// Deterministic partitioned reduction: no shared accumulator
int p = Environment.ProcessorCount;
double[] partial = new double[p];
Parallel.For(0, p, t => {
  int lo = (long)t * n / p;
  int hi = (long)(t + 1) * n / p;
  double s = 0.0;
  for (int i = lo; i < hi; i++) s += values[i] * values[i];
  partial[t] = s;
});
double total2 = partial.Sum();

// volatile: visibility/order constraint, NOT compound atomicity
private volatile bool _running = true;
\end{lstlisting}
\trap{For numerical reductions, prefer deterministic per-thread partials
or a lock-protected final merge. Avoid shared floating-point accumulators in
the loop body; their order is nondeterministic and usually slower.}
\trap{Beware false sharing: adjacent per-thread outputs can share a 64-byte
cache line and ping-pong between cores. Partition by contiguous output slices
or pad per-thread counters.}
\end{card}

\end{multicols}

% =====================================================================
\sectionbanner{File I/O, Strings, Regex \& JSON}
\begin{multicols}{2}

\begin{card}[acBlue]{bgBlue}{File \& Stream I/O}
\begin{lstlisting}[style=cs]
using System.IO;

// One-shot
string text   = File.ReadAllText("data.txt");
string[] lns  = File.ReadAllLines("data.txt");
byte[] bytes  = File.ReadAllBytes("data.bin");

// Async (preferred for I/O)
text = await File.ReadAllTextAsync("data.txt");
await File.WriteAllLinesAsync("results.txt", lns);

// Streaming for large files
using var sr = new StreamReader("big.csv");
string? line;
while ((line = sr.ReadLine()) is not null) Process(line);

using var sw = new StreamWriter("log.txt", append: false);
sw.WriteLine($"# {DateTime.UtcNow:O}");

// Binary I/O: raw IEEE 754
using var bw = new BinaryWriter(File.Create("vec.bin"));
bw.Write(n);                            // int header
for (int i = 0; i < n; i++) bw.Write(v[i]);

// Path / Directory
string dir = Path.Combine("results", "run01");
Directory.CreateDirectory(dir);
bool ex   = File.Exists("data.txt");
string fp = Path.GetFullPath("rel/path");
\end{lstlisting}
\note{For numerical data: prefer binary or text with \code{InvariantCulture}.
Never let locale settings affect persisted numeric files.}
\end{card}

\begin{card}[acGreen]{bgGreen}{Strings \& Number Formatting}
\begin{lstlisting}[style=cs]
string s = $"x = {x:F6}, n = {n:D6}";       // interpolation
string p = @"C:\data\file.txt";             // verbatim
string j = """{"x":3.14}""";                // raw

// Numeric format specifiers (most useful for science)
$"{v:G}"    // shortest round-trippable
$"{v:G15}"  // 15 sig digits  (good default for double)
$"{v:G17}"  // 17 sig digits  (full double precision)
$"{v:E6}"   // 1.234568E+003
$"{v:F4}"   // 1.2346
$"{i:D6}"   // zero-padded: 000042
$"{i:X8}"   // hex: 0000002A

// ALWAYS use InvariantCulture for persisted numeric data
double v = double.Parse(token, CultureInfo.InvariantCulture);

// Span-based parse (zero allocation, throws on invalid input)
ReadOnlySpan<char> tok = line.AsSpan(start, len);
double x = double.Parse(tok, NumberStyles.Float,
    CultureInfo.InvariantCulture);

// Formatting for persisted data
string txt = v.ToString("G17", CultureInfo.InvariantCulture);
writer.Write(txt);

// Multi-character search at SIMD speed
SearchValues<char> delims = SearchValues.Create(",;\t ");
int pos = line.IndexOfAny(delims);

// StringBuilder: avoids O(n^2) concatenation
var sb = new StringBuilder(capacity: 1024);
sb.Append("Hello").AppendLine()
  .AppendFormat(CultureInfo.InvariantCulture, "{0:F4}", x);
\end{lstlisting}
\note{\code{G17} for double is round-trip exact (\code{G9} for float).
\code{R} is the historical round-trip specifier; use \code{G17}
in new code.}
\end{card}

\begin{card}[acPurple]{bgPurple}{Regex --- Source-Generated}
\begin{lstlisting}[style=cs]
using System.Text.RegularExpressions;

// CLASSIC: pattern compiled at first use
var rx = new Regex(@"^\s*(-?\d+(?:\.\d+)?)\s*$",
                   RegexOptions.Compiled);
Match m = rx.Match(input);
if (m.Success) {
  double v = double.Parse(m.Groups[1].ValueSpan,
      CultureInfo.InvariantCulture);
}

// MODERN: source-generated; zero startup cost, AOT-safe
public partial class Parsers
{
  [GeneratedRegex(
    @"^\s*(?<x>-?\d+(?:\.\d+)?)\s+(?<y>-?\d+(?:\.\d+)?)\s*$",
    RegexOptions.IgnoreCase)]
  public static partial Regex Coord();
}
Match m2 = Parsers.Coord().Match(line);
if (m2.Success) {
  double x = double.Parse(m2.Groups["x"].ValueSpan,
      CultureInfo.InvariantCulture);
  double y = double.Parse(m2.Groups["y"].ValueSpan,
      CultureInfo.InvariantCulture);
}

foreach (Match hit in Parsers.Coord().Matches(text)) { }
string clean = Regex.Replace(s, @"\s+", " ");
\end{lstlisting}
\note{Use \code{Groups[i].ValueSpan} (not \code{Value}) to avoid
allocating a new string per match -- big win when parsing millions
of lines.}
\end{card}

\begin{card}[acTeal]{bgTeal}{System.Text.Json}
\begin{lstlisting}[style=cs]
using System.Text.Json;
using System.Text.Json.Serialization;

public sealed record Point(double X, double Y);
public sealed record Mesh(string Name,
                          IReadOnlyList<Point> Nodes);

string json = JsonSerializer.Serialize(mesh);
Mesh? back  = JsonSerializer.Deserialize<Mesh>(json);

// Options most relevant to scientific data
var opts = new JsonSerializerOptions {
  WriteIndented = true,
  PropertyNamingPolicy = JsonNamingPolicy.CamelCase,
  PropertyNameCaseInsensitive = true,
  // For strict interchange, also consider newer unmapped/duplicate policies.
  // CRUCIAL: round-trip NaN, Infinity safely
  NumberHandling = JsonNumberHandling.AllowNamedFloatingPointLiterals,
};
JsonSerializer.Serialize(mesh, opts);

// Source-generated serializer
// AOT-safe, zero reflection, ~2x faster
[JsonSerializable(typeof(Mesh))]
[JsonSerializable(typeof(Point))]
[JsonSourceGenerationOptions(WriteIndented = true)]
public partial class MeshCtx : JsonSerializerContext { }

string s = JsonSerializer.Serialize(mesh,
    MeshCtx.Default.Mesh);

// Stream-based for large files
await using var fs = File.Create("results.json");
await JsonSerializer.SerializeAsync(fs, mesh);
\end{lstlisting}
\note{Without \code{AllowNamedFloatingPointLiterals}, attempting to
serialize \code{double.NaN} throws \code{JsonException}.
Round-tripping FEM results almost always needs this.}
\end{card}

\begin{card}[acAmber]{bgAmber}{Dates, Times, Random \& Stopwatch}
\begin{lstlisting}[style=cs]
// PREFER DateTimeOffset over DateTime: explicit UTC offset
DateTimeOffset now = DateTimeOffset.UtcNow;
DateOnly d = DateOnly.FromDateTime(DateTime.Today);
TimeOnly t = TimeOnly.FromDateTime(DateTime.Now);

TimeSpan dt = now - start;
double secs = dt.TotalSeconds;

// ISO 8601 round-trip
string s = now.ToString("O", CultureInfo.InvariantCulture);
DateTimeOffset back = DateTimeOffset.Parse(s,
    CultureInfo.InvariantCulture);

// Stopwatch: high-res; NOT wall-clock
var sw = Stopwatch.StartNew();
DoWork();
sw.Stop();
double ms = sw.Elapsed.TotalMilliseconds;
// PORTABLE -- use TimeSpan, not raw ticks:
// Stopwatch.Frequency varies (1e7 on Win, ~1e9 on Linux);
// raw 'ticks' are not interchangeable with DateTime ticks.

// Random: thread-safe shared instance
double u = Random.Shared.NextDouble();   // [0, 1)
int    k = Random.Shared.Next(0, n);
// For reproducibility (e.g. Monte Carlo), use a SEEDED local
var rng = new Random(seed: 42);
// Crypto-strength: System.Security.Cryptography.RandomNumberGenerator
\end{lstlisting}
\trap{\code{Random.Shared} uses xoshiro -- excellent statistically but
NOT cryptographic and NOT reproducible across machines/runs. Always
seed a local \code{Random} for reproducible numerical experiments.}
\end{card}

\begin{card}[acRed]{bgRed}{Attributes, Reflection \& Source Gen}
\begin{lstlisting}[style=cs]
using System.Reflection;
using System.Diagnostics;

// CUSTOM ATTRIBUTE
[AttributeUsage(AttributeTargets.Property | AttributeTargets.Field)]
public sealed class UnitAttribute(string symbol) : Attribute
{
  public string Symbol { get; } = symbol;
}

public sealed class Reading {
  [Unit("m/s")] public double Velocity { get; init; }
  [Unit("Pa")]  public double Pressure { get; init; }
}

// REFLECTION: discover types, read attributes
foreach (PropertyInfo pi in typeof(Reading).GetProperties())
  if (pi.GetCustomAttribute<UnitAttribute>() is { } u)
    Console.WriteLine($"{pi.Name}: [{u.Symbol}]");

// Reflection invoke (slow; cache MethodInfo if hot)
MethodInfo mi = typeof(Math).GetMethod("Sqrt",
    [typeof(double)])!;
double r = (double)mi.Invoke(null, [4.0])!;   // 2.0

// SOURCE GENERATORS already in the BCL:
//   [GeneratedRegex]   regex without runtime compile
//   [JsonSerializable] reflection-free JSON
//   [LoggerMessage]    high-perf logging

// Activity: distributed-tracing spans (OpenTelemetry-compatible)
var src = new ActivitySource("MyApp");
using (var act = src.StartActivity("solve")) {
  act?.SetTag("n", n);
  Solve();
}
\end{lstlisting}
\note{Diagnose live processes with \code{dotnet-counters monitor},
\code{dotnet-trace collect}, and \code{dotnet-dump analyze}.
EventCounters expose GC stats, allocation rate, exceptions/sec.}
\end{card}

\end{multicols}

% =====================================================================
\sectionbanner{Numerical Best Practices \& Reference}
\begin{multicols}{2}

\begin{card}[acNavy]{bgNavy}{Version Scope}
\begin{itemize}
\item Scope: C\# 14 syntax and .NET 10 SDK/runtime/library surface relevant to numerical/scientific code.
\item Some APIs in \code{System.Numerics.Tensors} may be package-provided or preview-versioned depending on the exact SDK/package combination.
\item Treat every performance rule as a hypothesis until BenchmarkDotNet or an application-level benchmark confirms it on the target CPU, OS, and runtime.
\item For public libraries, keep nullable analysis enabled, document ownership of buffers, and separate safe APIs from unsafe fast paths.
\end{itemize}
\end{card}

\begin{card}[acAmber]{bgAmber}{Numerical Best Practices}
\begin{itemize}
\item Default to \code{double}. Use \code{float} or \code{Half}
  only when memory bandwidth dominates and you have measured the win.
\item \code{decimal} is base-10 and slow -- use it for money,
  not science.
\item Always pass \code{CultureInfo.InvariantCulture} when parsing or
  formatting persisted numeric data.
\item Use compensated (Kahan or Neumaier) or pairwise summation for
  long sums; naive sum's error grows like $O(n \cdot \epsilon)$.
\item Use \code{T.FusedMultiplyAdd(a, b, c)} or \code{Math.FusedMultiplyAdd}
  when one-rounding semantics improve robustness; measure throughput on target hardware.
\item Reach for \code{TensorPrimitives} before writing manual SIMD: it handles
  tails and uses optimized implementations when available.
\item Pre-size collections; \code{List<T>} resize copies are silent killers.
\item Inside hot loops: indexed \code{for}, \code{ReadOnlySpan<T>}
  parameters, no closures, no \code{string.Split}.
\item Pin allocations \emph{outside} loops; rent from \code{ArrayPool}
  for large transients; use native aligned memory when exact SIMD alignment
  is required for interop or kernels.
\item \code{readonly struct} for tiny vectors -- prevents costly
  defensive copies at \code{in} parameter sites.
\item Use \code{checked\{\}} where overflow indicates a bug;
  \code{unchecked\{\}} only when wraparound is intended.
\item Always pass \code{CancellationToken} through async chains.
\item Validate inputs once at the public API boundary; trust the
  kernel.
\end{itemize}
\sub{Measurement discipline}
\begin{itemize}
\item Release build, no debugger, warm-up before timing.
\item BenchmarkDotNet for microbenchmarks; \code{Stopwatch} for coarse.
\item Track allocations as a first-class metric -- \code{dotnet-counters}.
\item Validate correctness BEFORE optimising. Fast wrong code is worthless.
\end{itemize}
\end{card}

\begin{card}[acBlue]{bgBlue}{Key Interfaces \& Selection Rules}
\fontsize{5.95pt}{7.6pt}\selectfont
\begin{tabular}{@{}p{27mm}@{\;}p{26mm}@{}}
\rowcolor{bgHdr}\textbf{Interface} & \textbf{Use}\\
\code{IEquatable<T>}        & typed Equals; pair with \code{GetHashCode}\\
\code{IComparable<T>}        & natural total order; enables \code{Sort}\\
\code{IComparer<T>}          & external comparator\\
\code{IEqualityComparer<T>} & key equality for dict/hash\\
\code{IEnumerable<T>}       & lazy iteration; supports \code{foreach}\\
\code{IReadOnlyList<T>}     & indexed, immutable public APIs\\
\code{IList<T>}             & only when caller must mutate\\
\code{IDisposable}          & deterministic release; pair with \code{using}\\
\code{IAsyncDisposable}     & async cleanup; \code{await using}\\
\code{IAsyncEnumerable<T>}  & \code{await foreach} streaming\\
\code{INumberBase<T>}       & arithmetic, \code{Create*} conversions\\
\code{INumber<T>}           & adds comparison + sign\\
\code{IFloatingPointIeee754<T>} & adds Sqrt/Sin/Pi/NaN/FMA\\
\code{ISpanFormattable}     & zero-alloc \code{TryFormat}\\
\code{ISpanParsable<T>}     & zero-alloc \code{Parse} from span\\
\end{tabular}
\sub{Selection rule}
{\fontsize{5.7pt}{7pt}\selectfont
Accept the WEAKEST interface that satisfies the caller. For numerical
kernels prefer \code{ReadOnlySpan<T>} or \code{Span<T>} over any
interface -- this bypasses virtual dispatch entirely and enables
implicit \code{T[]} conversions.}
\sub{INumberBase Create methods}
{\fontsize{5.7pt}{7pt}\selectfont
\code{T.CreateChecked<U>(x)} throws \code{OverflowException} outside range.\\
\code{T.CreateSaturating<U>(x)} clamps to \code{T.MinValue} / \code{T.MaxValue}.\\
\code{T.CreateTruncating<U>(x)} bit-truncates / wraps -- use when you
explicitly want modular arithmetic.}
\end{card}

\begin{card}[acPurple]{bgPurple}{Numerical Correctness Traps}
\begin{itemize}
\item \code{x == double.NaN} is always \code{false} -- use \code{double.IsNaN(x)}.
\item Sorting with NaN: \code{Array.Sort} uses \code{IComparable},
  which puts NaN \emph{first}. \code{<} returns \code{false} for any
  comparison with NaN. The two disagree.
\item \code{0.1 + 0.2 != 0.3} -- never compare floats with \code{==}
  for equality; use a tolerance: \code{Math.Abs(a - b) < tol * Math.Max(Math.Abs(a), Math.Abs(b))}.
\item \code{Math.Round} uses \emph{banker's rounding} by default
  (round-to-even). For school rounding pass
  \code{MidpointRounding.AwayFromZero}.
\item Subtraction of nearly-equal numbers loses precision
  (\emph{catastrophic cancellation}). Reformulate:
  $\sqrt{x^2 + y^2}$ as \code{double.Hypot(x, y)};
  $1 - \cos x$ as $2\sin^2(x/2)$.
\item For sums of many terms: prefer pairwise or
  Kahan summation. error bounds depend on conditioning, but naive summation can grow like
  $O(n\,\epsilon)$ while pairwise and compensated schemes usually reduce
  accumulation error substantially.
\item \code{Math.FusedMultiplyAdd(a, b, c)} is specified as one rounding;
  \code{a * b + c} is normally rounded after multiply and again after add.
\item Integer overflow is silent by default. Wrap risky arithmetic
  in \code{checked\{\}} or compile with \code{<CheckForOverflowUnderflow>}.
\item \code{double} keys in \code{Dictionary}: \code{double.NaN.Equals(double.NaN)}
  is \code{true}, while \code{double.NaN == double.NaN} is \code{false}.
  Avoid floating-point keys unless you define the exact equality semantics.
\item Persisted \code{double} round-trips need \code{G17}, not \code{G15}.
  \code{float} needs \code{G9}.
\end{itemize}
\end{card}

\end{multicols}

\end{document}
